\section{Limits of Coherence-Based Gating}
\label{sec:limits}

Section~\ref{sec:gateaware} showed an adversary whose optimum coherence is below
chance. This section argues that no amount of threshold engineering recovers
from that, and identifies what a successor defence would have to bring instead.

\subsection{Proposition 1, informally}

Fix a question with a closed answer space of size $C$. AIP's gate is a decision
rule computed from a channel's per-question coherence $q$, the probability that
two agents on the channel, both wrong, are wrong the same way. Honest agents
failing independently collide at the chance rate $q_{\mathrm{ind}} = 1/(C-1)$.

\begin{quote}
\emph{Let $\mathcal{A}$ be an adversarial channel with $q_{\mathcal{A}} \leq
q_{\mathrm{ind}}$, and $\mathcal{H}$ a population of honest agents failing
independently on the same space. Then for any decision rule computed solely from
per-question coherence, the sum of its false-positive and false-negative rates is
at least $1 - d_{TV}$, where $d_{TV}$ is the total variation distance between the
statistic's distributions under the two hypotheses, and $d_{TV} \to 0$ as
$q_{\mathcal{A}} \to q_{\mathrm{ind}}$.}
\end{quote}

The argument is short. Collision counts among jointly-dissenting peers are a
sufficient statistic for $q$ under the exchangeability the estimator already
assumes; by the Neyman--Pearson lemma the best achievable separation between two
hypotheses on a common statistic is governed by the total variation distance
between its distributions; the count is Binomial$(n, q)$, so $d_{TV}$ is
continuous in $q$ and vanishes as the two parameters meet. The bound therefore
constrains \emph{every} rule measurable with respect to that statistic, not the
particular threshold we chose. The formal setup, the exchangeability condition in
full and the finite-$n$ statement are in the artifact's technical note; the
substance is that $d_{TV}$ degrades as $\sqrt{n}$ in the number of
jointly-dissenting observations, which is bounded by the swarm size --- so the
escape route is longitudinal rather than per-question --- and that at
$q_{\mathcal{A}} = q_{\mathrm{ind}}$ exactly, $d_{TV}=0$ for every $n$.

This is why M1 helps a little and M2 not at all: soft gating recovers accuracy in
the band below the \emph{calibrated ceiling} but still above chance, where the
statistic carries signal, while at and below chance there is nothing for any
function of $q$ to find and randomising the threshold only adds variance.

\begin{figure}[t]
\centering
\includegraphics[width=\columnwidth]{evasion_band.pdf}
\caption{The evasion band at the anchor, $f = \numEvasionAnchorF$. Markers are
the measured grid; nothing is interpolated between them. The dashed line is the
calibrated honest ceiling the gate fires against, the dotted line is the chance
coherence $1/(C-1)$ of independent wrong answers, and the rug shows the
per-pair honest coherence with underpowered pairs ($n_{\text{both wrong}} < 30$,
rule R2) drawn in grey. On MedQA the attacker's optimum sits below both lines;
on MATH-500 the gate is near-inert and the attacker's optimum is a fully
coherent channel, so there is nothing there to evade.}
\label{fig:evasion}
\end{figure}

\subsection{What the measurement can and cannot say}

The anchor is a single grid point and we are explicit about its resolution: the
parameterisation reaches only $\numEvasionQDistinct$ distinct coherence values
inside $[\numEvasionBandLo, \numEvasionBandHi]$ and cannot express a channel less
coherent than $\numEvasionQFloorMc$, the interior minimum of $q(p)$. The optimum
is \emph{measured} at a grid point, not interpolated; we claim nothing about its
shape between points, and whether a differently constructed adversary could reach
lower $q$ is untested.

\subsection{Inversion and evasion are one law, read from two answer spaces}

The mechanism's success and its defeat are not two findings. They are the same
inequality evaluated on two answer spaces.

The inequality is whether the adversary's coherence exceeds the honest
receiver-conditioned coincidence on that space. On open-numeric answers under
negation it does --- $\hat q = \numGsmkNegQhat$ against an honest floor of
$\numGsmkHonestFloor$ --- so there is a separable channel, the gate fires, and
inversion converts the adversary into a contributor. Every benchmark where AIP
separates from the discard family is on this side of the line. On four-way
multiple choice it does not: $\numQhatMcMean$ against $\numMmluHonestFloor$. The
adversary is \emph{more} coordinated in absolute terms than on open text and
still carries no channel, because the honest agents are more coordinated too.
\textbf{[G2']}

The gate-aware adversary is what happens when an adversary chooses the second
side of that line deliberately: its optimum $\numEvasionQMedqa$ sits below the
honest floor and below chance $\numEvasionChance$. It is not evading a threshold,
it is moving into the regime where no threshold on this statistic has anything to
work with --- exactly the regime Proposition~1 describes. So the practitioner's
question is not ``does inversion work'' but ``does this answer space put my
adversary above the honest floor''; on a binary space the answer is no by
construction, which is why majority vote beats AIP on BoolQ. The same law
predicts the win and the loss. \textbf{[D3]}

\subsection{Two escape classes, both untested}

The proposition constrains defences whose statistic is per-question channel
coherence. It says nothing about defences drawing on information that statistic
does not contain. Two such classes are identifiable, and both are stated here as
\textbf{hypotheses that have not been evaluated against the gate-aware
adversary}.

\textbf{External information.} Identity verification, stake, or a ground-truth
anchor breaks the symmetry the proposition relies on by supplying evidence from
outside the channel. This is precisely the gap AgentChain~\cite{agentchain}
names as unsolved when it admits that its guarantees stop at $f \geq N/2$
(Section~\ref{sec:related}). Our contribution here is not to close that gap but
to give a reason to think a coherence-based defence cannot: the information
required is not present in the statistic, so no better estimator of the statistic
will supply it.

\textbf{Longitudinal statistics.} Cross-task reputation accumulated over rounds
is not a function of per-question coherence and is therefore outside the
proposition's scope. Our own sleeper experiment is suggestive in this direction
and against our prior expectation: reputation decay degraded least of the
memory-carrying defenders, at $\numSleeperDropReputation$ against AIP-gated's
$\numSleeperDropAip$, a spread of $\numSleeperDropSpread$ that only just clears
the $\numFloorWorst$ floor. That is a marginal result pointing at an untested
hypothesis, and it is worth exactly that much. A longitudinal defence has not
been run against the gate-aware adversary and we do not claim it survives.

\subsection{What this means for the mechanism}

Sections~\ref{sec:evaluation}--\ref{sec:limits} describe a defence with a
characterised operating envelope rather than a solution: close to free when idle,
large when the adversary is coherent and the answer space open, useless when the
space is binary, and defeated by an adversary that reads its calibration constant
and sits below it. That last is a property of the statistic, not a bug to be
patched.
